%%% Hlavní soubor. Zde se definují základní parametry a odkazuje se na ostatní části. %%%

%% Verze pro jednostranný tisk:
% Okraje: levý 40mm, pravý 25mm, horní a dolní 25mm
% (ale pozor, LaTeX si sám přidává 1in)
\documentclass[12pt,a4paper]{report}
\setlength\textwidth{145mm}
\setlength\textheight{247mm}
\setlength\oddsidemargin{15mm}
\setlength\evensidemargin{15mm}
\setlength\topmargin{0mm}
\setlength\headsep{0mm}
\setlength\headheight{0mm}
% \openright zařídí, aby následující text začínal na pravé straně knihy
\let\openright=\clearpage

%% Pokud tiskneme oboustranně:
% \documentclass[12pt,a4paper,twoside,openright]{report}
% \setlength\textwidth{145mm}
% \setlength\textheight{247mm}
% \setlength\oddsidemargin{15mm}
% \setlength\evensidemargin{0mm}
% \setlength\topmargin{0mm}
% \setlength\headsep{0mm}
% \setlength\headheight{0mm}
% \let\openright=\cleardoublepage

%% Použité kódování znaků: obvykle latin2, cp1250 nebo utf8:
\usepackage[utf8]{inputenc}

%% Ostatní balíčky
\usepackage{graphicx}
\usepackage{amsthm}

%% Balíček hyperref, kterým jdou vyrábět klikací odkazy v PDF,
%% ale hlavně ho používáme k uložení metadat do PDF (včetně obsahu).
%% POZOR, nezapomeňte vyplnit jméno práce a autora.
\usepackage[ps2pdf,unicode]{hyperref}   % Musí být za všemi ostatními balíčky
\hypersetup{pdftitle=Název práce}
\hypersetup{pdfauthor=Jméno Příjmení}

%%% Drobné úpravy stylu

% Tato makra přesvědčují mírně ošklivým trikem LaTeX, aby hlavičky kapitol
% sázel příčetněji a nevynechával nad nimi spoustu místa. Směle ignorujte.
\makeatletter
\def\@makechapterhead#1{
  {\parindent \z@ \raggedright \normalfont
   \Huge\bfseries \thechapter. #1
   \par\nobreak
   \vskip 20\p@
}}
\def\@makeschapterhead#1{
  {\parindent \z@ \raggedright \normalfont
   \Huge\bfseries #1
   \par\nobreak
   \vskip 20\p@
}}
\makeatother

% Toto makro definuje kapitolu, která není očíslovaná, ale je uvedena v obsahu.
\def\chapwithtoc#1{
\chapter*{#1}
\addcontentsline{toc}{chapter}{#1}
}

\begin{document}

% Trochu volnější nastavení dělení slov, než je default.
\lefthyphenmin=2
\righthyphenmin=2

%%% Titulní strana práce

\pagestyle{empty}
\begin{center}

\large

Charles University in Prague

\medskip

Faculty of Mathematics and Physics

\vfill

{\bf\Large BACHELOR THESIS}

\vfill

\centerline{\mbox{\includegraphics[width=60mm]{./img/logo.eps}}}

\vfill
\vspace{5mm}

{\LARGE Daniel Lovásko}

\vspace{15mm}

% Název práce přesně podle zadání
{\LARGE\bfseries CTF implementation and use in the FreeBSD operating system}

\vfill

% Název katedry nebo ústavu, kde byla práce oficiálně zadána
% (dle Organizační struktury MFF UK)
Department of Applied Mathematics

\vfill

\begin{tabular}{rl}

Supervisor of the bachelor thesis: & Mgr. Petr Baudiš \\
\noalign{\vspace{2mm}}
Study programme: & Informatics \\
\noalign{\vspace{2mm}}
Specialization: & Programming \\
\end{tabular}

\vfill

% Zde doplňte rok
Prague 2015

\end{center}

\newpage

%%% Následuje vevázaný list -- kopie podepsaného "Zadání bakalářské práce".
%%% Toto zadání NENÍ součástí elektronické verze práce, nescanovat.

%%% Na tomto místě mohou být napsána případná poděkování (vedoucímu práce,
%%% konzultantovi, tomu, kdo zapůjčil software, literaturu apod.)

\openright

\noindent
Dedication.
%\begin{itemize}
%\item[Petr Baudis]
%\item[George Neville-Neil]
%\item[Dominik Lukac]
%\item[Matej Vitasek]
%\item[Ladislav Papay]
%\item[Pedro Giffuni]
%\item[Robert Mustacchi]
%\item[Robert Watson]
%\end{itemize}

\newpage

%%% Strana s čestným prohlášením k bakalářské práci

\vglue 0pt plus 1fill

\noindent
I declare that I carried out this bachelor thesis independently, and only with the cited
sources, literature and other professional sources.

\medskip\noindent
I understand that my work relates to the rights and obligations under the Act No.
121/2000 Coll., the Copyright Act, as amended, in particular the fact that the Charles
University in Prague has the right to conclude a license agreement on the use of this
work as a school work pursuant to Section 60 paragraph 1 of the Copyright Act.

\vspace{10mm}

\hbox{\hbox to 0.5\hsize{%
In ........ date ............
\hss}\hbox to 0.5\hsize{%
signature of the author
\hss}}

\vspace{20mm}
\newpage

%%% Povinná informační strana bakalářské práce

\vbox to 0.5\vsize{
\setlength\parindent{0mm}
\setlength\parskip{5mm}

Název práce:
Název práce
% přesně dle zadání

Autor:
Jméno a příjmení autora

Katedra:  % Případně Ústav:
Název katedry či ústavu, kde byla práce oficiálně zadána
% dle Organizační struktury MFF UK

Vedoucí bakalářské práce:
Jméno a příjmení s tituly, pracoviště
% dle Organizační struktury MFF UK, případně plný název pracoviště mimo MFF UK

Abstrakt:
% abstrakt v rozsahu 80-200 slov; nejedná se však o opis zadání bakalářské práce

Klíčová slova:
% 3 až 5 klíčových slov

\vss}\nobreak\vbox to 0.49\vsize{
\setlength\parindent{0mm}
\setlength\parskip{5mm}

Title:
% přesný překlad názvu práce v angličtině

Author:
Jméno a příjmení autora

Department:
Název katedry či ústavu, kde byla práce oficiálně zadána
% dle Organizační struktury MFF UK v angličtině

Supervisor:
Jméno a příjmení s tituly, pracoviště
% dle Organizační struktury MFF UK, případně plný název pracoviště
% mimo MFF UK v angličtině

Abstract:
% abstrakt v rozsahu 80-200 slov v angličtině; nejedná se však o překlad
% zadání bakalářské práce

Keywords:
% 3 až 5 klíčových slov v angličtině

\vss}

\newpage

%%% Strana s automaticky generovaným obsahem bakalářské práce. U matematických
%%% prací je přípustné, aby seznam tabulek a zkratek, existují-li, byl umístěn
%%% na začátku práce, místo na jejím konci.

\openright
\pagestyle{plain}
\setcounter{page}{1}
\tableofcontents

%%% Jednotlivé kapitoly práce jsou pro přehlednost uloženy v samostatných souborech
\chapter*{Introduction}
\addcontentsline{toc}{chapter}{Introduction}


\chapter{Compact C Type Format}

\section{History and motivation}
\section{Format description}
\subsection{ELF section}
\subsection{Header}
\subsection{Labels and uniquification}
\subsection{Function objects}
\subsection{Data objects}
\subsection{Types}
\subsection{Variable data}
\subsubsection{Integer number}
\subsubsection{Floating point number}
\subsubsection{Enumeration}
\subsubsection{Structure and union}
\subsubsection{Function}
\subsection{Strings}
\section{Contemporary usage}


\chapter{libctf}

\section{Implementation}
\section{Infrastructure}
\subsection{Version control}
\subsection{Build system}
\subsection{Documentation system}
\section{Licence}
\section{Performance}
\section{Memory usage}
\subsection{Memory leak prevention}
\section{Correctness}
\subsection{Unit tests}
\subsection{Code coverage}


\chapter{Toolset}

\section{ctfdump}
\section{ctfstats}
\section{ctfconvert}
\section{ctfmerge}


\chapter{Pretty-printing kernel data structures}
Memory correctness plays the key role in any algorithm. Being able to swiftly
and efficiently examine arbitrarily complex data structures is a mandatory
trait of every modern toolchain. The GNU Debugger {\tt gdb}, Modular Debugger
{\tt mdb} and the LLVM Debugger {\tt lldb} fulfil the expectations with respect
to userland processes. Unfortunately, the kernel debugger DDB on FreeBSD (and
other systems using it) has been shipped without such functionality from
its inception many decades ago. We aim to solve this issue by creating a
mechanism to pretty-print data structures contained in the currently loaded
kernel image.

\section{Approach comparison}
Currently, DDB offers two techniques to access the memory which are affected by
many limitations, either lack of modularity or excessive straightforwardness.

\subsection{Low-level examination}
The {\tt examine} command [referencia!] in conjecture with its options such as
{\tt x} or {\tt f} is used to examine optional number of memory blocks. The
fact that the user must specify the encoding of the memory is a .

This feature of DDB is undoubtedly important in scenarios when a specific
bit or byte needs to be analysed, but falls short in examining complex
structures, as it literally forces user to use pen and paper to map the output
values to the data structure.

\subsection{Hard-coded structures}
The {\tt show} command [referencia!] improves upon the simplicity of th {\tt
examine} command by providing support for pretty-printing a restricted set of
data structures. Such structures are:
\begin{itemize}
	\item bio
	\item mutex
\end{itemize}

The problem with this approach lies in its {\it code change linearity}: every
time a data structures member changes, the corresponding code in DDB needs to
adapt. The same applies in situations when a new data structure is introduced -
a new pretty-printing code needs to be written.

\subsection{Universal CTF solution}
The logical extension of the {\tt show} command is to provide a way of encoding
type information during the compilation process and being able to read such
data later during the debugger runtime - a task that CTF is a great fit for -
so that there will be no need to alter the DDB source code in case of a
structure change in the kernel code-base. We therefore implemented such
approach by using the {\tt libctf} described in the Chapter 2. Following
sections describe and discuss the implementation details.

\section{Kernel-space adaptation}
\section{Type lookup}
\section{Pretty printing}
\section{Data structure speculation}
\section{Views}
\section{Corrupted structures}



% Ukázka použití některých konstrukcí LateXu (odkomentujte, chcete-li)
% \include{example}

\chapter*{Conclusion}
\addcontentsline{toc}{chapter}{Conclusion}


%%% Seznam použité literatury
%%% Seznam použité literatury je zpracován podle platných standardů. Povinnou citační
%%% normou pro bakalářskou práci je ISO 690. Jména časopisů lze uvádět zkráceně, ale jen
%%% v kodifikované podobě. Všechny použité zdroje a prameny musí být řádně citovány.

\def\bibname{Bibliography}
\begin{thebibliography}{99}
\addcontentsline{toc}{chapter}{\bibname}

%\bibitem{lamport94}
%  {\sc Lamport,} Leslie.
%  \emph{\LaTeX: A Document Preparation System}.
%  2. vydání.
%  Massachusetts: Addison Wesley, 1994.
%  ISBN 0-201-52983-1.

\end{thebibliography}


%%% Tabulky v bakalářské práci, existují-li.
\chapwithtoc{List of Tables}

%%% Použité zkratky v bakalářské práci, existují-li, včetně jejich vysvětlení.
\chapwithtoc{List of Abbreviations}

%%% Přílohy k bakalářské práci, existují-li (různé dodatky jako výpisy programů,
%%% diagramy apod.). Každá příloha musí být alespoň jednou odkazována z vlastního
%%% textu práce. Přílohy se číslují.
\chapwithtoc{Attachments}

\openright
\end{document}
